\section{Introduction}

Repeated reads and state-changing writes often meet in the same backend. A cache can serve some reads without sending their queries to the database. For a point-of-sale (POS) system, the saved database work is useful only if the application understands how old a cached answer can be. A terminal may show an item as available after another terminal has sold the last portion. The order still has to pass an authoritative inventory check. In the meantime, a customer can make an order attempt based on an item that has already sold out.

We study this trade-off in an Express POS backend that stores menu recipes, ingredient stock, and orders in PostgreSQL. Several items can use the same ingredients, so one order can affect the cached availability of related items. The baseline reads through PostgreSQL. The other policies use Redis in a cache-aside design. A cache-aside reader checks Redis first and queries PostgreSQL on a miss. The three cached policies use expiration alone, expiration with post-commit deletion, or expiration with deletion and a generation check that protects against delayed refills.

The study asks three questions. First, can a read that overlaps an order repopulate Redis with obsolete stock after the order has acknowledged, and does a generation check prevent that refill in the same schedule? Second, how many correct stockout rejections follow an outdated display in a controlled last-portion scenario? Third, how do the policies change endpoint latency, request delivery, and database query activity at matched arrival rates, and how sensitive are two of those modes to TTL? A screen identifies the rate region where delivery shortfalls appear. Repeated trials characterize variation there and at a lower reference rate.

The contribution is a reproducible application-level comparison, including a replayable interleaving and a way to count rejected checkout attempts attributable to false availability in the experiment. Generation checks and conditional writes build on prior cache consistency mechanisms \cite{nishtala2013}. The implementation tests how those mechanisms work with shared ingredient dependencies in this POS backend. The requests and dataset reflect actual trial data. The measured services and workload generator ran on one Linux host as local processes.
